% from ega4-16-fr.tex
\oldpage[IV-4]{35}
\item[(iii)] Pour tout ouvert $U$ de $X$, tout couple d'entiers $p$, $q$ et
tout couple de sections
$\omega'_p\in\Gamma(U,\Omega_{X/S}^p)$,
$\omega''_q\in\Gamma(U,\Omega_{X/S}^q)$, on a
\[
  d(\omega'_p\wedge\omega''_q)
  =(d\omega'_p)\wedge\omega''_q
  +(-1)^p\omega'_p\wedge d\omega''_q.
  \tag{16.6.2.1}
  \label{IV.16.6.2.1}
\]
\end{enumerate}
En outre, $d$ est un endomorphisme de $\psi^*(\mathcal{O}_X)$-Modules
gradués, de degré $+1$.
\end{theorem}

Supposons prouvée l'existence de l'endomorphisme $d$. Pour tout ouvert affine
$U$ de $X$, toute section de $\Omega_{X/S}^p$ au-dessus de $U$ est (en vertu
de (i)) combinaison linéaire d'un nombre fini d'éléments de la forme
$g(df_1\wedge df_2\wedge\cdots\wedge df_p)$, où $g$ et les $f_i$ sont des
sections de $\mathcal{O}_X$ au-dessus de $U$ (0, 20.4.7). Les conditions (i)
et (iii) montrent alors, par récurrence sur $p$, que l'on a nécessairement
\[
  d(g(df_1\wedge df_2\wedge\cdots\wedge df_p))
  =dg\wedge df_1\wedge df_2\wedge\cdots\wedge df_p.
  \tag{16.6.2.2}
  \label{IV.16.6.2.2}
\]
Cela prouve donc l'\emph{unicité} de $d$ et la dernière assertion du théorème.
En vertu de cette propriété d'unicité, pour démontrer l'existence de $d$, on
peut se borner au cas où $S=\operatorname{Spec}(A)$ et
$X=\operatorname{Spec}(B)$ sont affines. Or (Bourbaki, \emph{Alg.}, chap.
III, 3\textsuperscript{e} éd., \S~10) pour définir une $A$-antidérivation
$D$ de degré $+1$ d'une algèbre extérieure $\bigwedge(M)$ (où $M$ est un
$B$-module et $B$ une $A$-algèbre), cette antidérivation prenant ses valeurs
dans une $A$-algèbre \emph{graduée anticommutative}
$C=\bigoplus_{n=0}^{\infty}C_n$, dont les éléments de degré 1 sont de carré
nul, il suffit de se donner arbitrairement une $A$-dérivation $D_0$ de $B$
dans $C_1$ et un $A$-homomorphisme $D_1$ de $M$ dans $C_2$; il existe alors
une $A$-antidérivation $D$ et une seule de $\bigwedge(M)$ dans $C$
coïncidant avec $D_0$ dans $B$ et avec $D_1$ dans $M$.

Dans le cas actuel, $D_0$ est nécessairement égal à $d_{B/A}$ en vertu de
(ii); tout revient à voir, compte tenu de (16.6.2.2), qu'il y a un
$A$-homomorphisme $u$ de $\Omega_{B/A}^1$ dans $\Omega_{B/A}^2$ tel que
l'on ait
\[
  u(g.df)=dg\wedge df
  \tag{16.6.2.3}
  \label{IV.16.6.2.3}
\]
quels que soient $f$, $g$ dans $A$; il suffira pour cela de montrer qu'il
existe un $A$-homomorphisme
$v:B\otimes_A\Omega_{B/A}^1\to\Omega_{B/A}^2$ tel que
\[
  v(g.\omega)=dg\wedge\omega
  \tag{16.6.2.4}
  \label{IV.16.6.2.4}
\]
pour $g\in B$ et $\omega\in\Omega_{B/A}^1$. Finalement, comme
$\Omega_{B/A}^1=\mathfrak{I}/\mathfrak{I}^2$ (où
$\mathfrak{I}=\mathfrak{I}_{B/A}$ est le noyau de l'homomorphisme canonique
$B\otimes_A B\to B$) et que $\Omega_{B/A}^1$ est engendré par les éléments
de la forme $g.df$, il suffit de définir un $A$-homomorphisme
$w:B\otimes_A(B\otimes_A B)\to\Omega_{B/A}^2$ tel que
\[
  w(g'\otimes g\otimes f)=dg'\wedge(g.df)
  \tag{16.6.2.5}
  \label{IV.16.6.2.5}
\]
et tel que $w$ soit nul dans l'image de $B\otimes_A\mathfrak{I}^2$. Or,
comme le second membre de (16.6.2.5) est $A$-trilinéaire en $g'$, $g$, $f$,
l'existence de $w$ vérifiant (16.6.2.5) est immédiate. Comme, d'autre part,
$\mathfrak{I}$ est engendré par les éléments $1\otimes x-x\otimes1$
($x\in B$), on est ramené à vérifier que lorsque
$z=(1\otimes x-x\otimes1)(1\otimes y-y\otimes1)$, on a
$w(g'\otimes z)=0$. Or, comme
$z=1\otimes(xy)+(xy)\otimes1-x\otimes y-y\otimes x$, la formule (16.6.2.4)
montre qu'il
